\documentclass[12pt]{article}

\title{AlgoBOWL Software Design}
\date{24 April 2017}
\author{Jack Rosenthal}

\usepackage{fontspec}
\usepackage{hyperref}
\usepackage[margin=25mm]{geometry}
\usepackage{wasysym}
\usepackage{float}
\usepackage{graphicx}

\newcommand\ctrltitle[1]{\par\medskip\texttt{\large #1}\\}
\newcommand\ctrlauth[1]{Authorization: #1\\}
\newcommand\ctrldesc[1]{$\rhd$ \textsl{#1}\par\medskip}

\setlength\parindent{0pt}
\setlength\parskip{5pt plus 1pt minus 1pt}

\begin{document}

\maketitle

\section{Introduction}

AlgoBOWL is a web application used in CSCI-406: Algorithms. Students are given
a NP-hard problem to solve by the professor and work in small groups (typically
of three students) to create and implement a heuristic to get the best solution
they can. Once they have done this, each of the groups:
\begin{enumerate}
    \item Designs an input for the other groups to run and uploads it to the
        AlgoBOWL webpage
    \item Downloads all the inputs from the other groups and runs it on their
        program
    \item Uploads their outputs for the corresponding inputs
    \item Verifies each of the outputs for their input and marks it as either
        \emph{Accepted} or \emph{Rejected}
    \item Open verification is enabled, and anyone can view all inputs and
        outputs.
\end{enumerate}

To keep the process flowing, the students are given deadlines they need to
complete each of these steps by.

At the end of the assignment, each of the students are given an opportunity to
rank their fellow team members' efforts as a percentage. The instructor then
downloads these effort percentages and assigns grades to each student.

\section{Database Model Design}

Figure \ref{erd} shows a (Peter Chen style) entity relationship diagram showing
the different database models that should be used to implement the software.

\begin{figure}
    \includegraphics[width=\textwidth]{graphics/algo_erd.pdf}
    \caption{The database model}
    \label{erd}
\end{figure}

Notice the following:

\begin{itemize}
    \item The attribute \texttt{group\_id} of the \texttt{User} model specifies
        the group ID that the student is a participant of. The user can be a
        part of multiple groups if they fail the course and have to take it
        again (or they are Mehta and participate every semester).
    \item The \texttt{Group} model relates the numerical group ID to the team's
        name. Only the instructor may have means of discovering both sides of
        this relation as the names of the students in each group are not
        anonymous but their results must be. Each group is related to one
        \texttt{Competition}.
    \item Both the \texttt{Input} and \texttt{Output} models store the files
        using Depot.
    \item The \texttt{Output} model contains both a \texttt{score} attribute
        and a \texttt{data} attribute. When an output is uploaded, the first
        line is separated from the remainder of the file and parsed as a
        numerical value. This is explained further later on in the controller
        logic.
    \item The primary key if the \texttt{Output} model is the combination of
        both the group that uploaded the output and the group that the input
        corresponds to. There is no \texttt{id} attribute.
    \item The \texttt{verification} attribute of the \texttt{Output} model is a
        enumerable type; being one of \emph{Waiting}, \emph{Accepted}, or
        \emph{Rejected}.
    \item The evaluation is a recursive relation between users and other users,
        with data associated (eval percentage). Not shown in the ERD is that
        there is a constraint that the sum of evaluations from one user to N
        others in the same group is equal to 100.
\end{itemize}

\clearpage

\section{Controller Design}

\ctrltitle{GET /}
\ctrlauth{Anyone}
\ctrldesc{Generate the rankings page}

When the landing page for AlgoBOWL is loaded, a table of rankings for each
current competition is shown to the user. This rankings table is accessible
\emph{regardless of weather the current user is a CSCI-406 student} to allow
others to look in at the progress of AlgoBOWL.

Routes should also be added for specific competition ranking pages.

An example table is shown in Figure \ref{algotbl}. In this table:
\begin{itemize}
    \item A ranking of \textsf{N} indicates that the output was not submitted yet
    \item A ranking of \textsf{R} indicates that the output was rejected
    \item A ranking of a number with a \textsf{(W)} next to it indicates the output is waiting for acceptance
    \item A ranking of a number without a \textsf{(W)} next to it indicates the output was accepted
    \item Groups highlighted in red have at least one rejected output
    \item Groups highlighted in orange have at least one output they have not submitted yet
    \item Groups highlighted in green have all their outputs accepted
\end{itemize}

\begin{figure}
    \includegraphics[width=\textwidth]{graphics/algotbl}
    \caption{An AlgoBOWL Rankings Table}
    \label{algotbl}
\end{figure}

A few important notes on how this table was generated:

\begin{itemize}
    \item For a particular input, if any teams tie, they should share the same
        rank and the ranks of all the teams should offset accounting for the
        number of tied ranks. For example, if two teams tie for first place,
        the next place is third place, not second. In other words, if $N$ teams
        did better on a particular input, then the rank of that team for the
        input is defined to be $N + 1$.
    \item When a team has rejected outputs or has not uploaded all of their
        outputs, they cannot receive a finite value for sum of ranks, because
        they have not yet completed the competition.
    \item When a team has not had their output accepted yet, they can be given
        a rank compared to the others, but the table should indicate it is not
        accepted yet.
    \item When two teams tie for their sum of ranks, a reasonable secondary
        sorting key should be used to show their order in the table. This
        secondary key (sorted by minimum) is calculated as $R - 2A - W$, where $R$
        is the number of rejected outputs for that team, $A$ is the number of
        accepted outputs for that team, and $W$ is the number of outputs
        waiting acceptance for that team. In the case that this secondary sort
        key is equivalent as well, the group's number is used as the third key.
    \item Group numbers should not be shown next to their scores (but instead
        their group name), as this preserves student anonymity.
\end{itemize}

\ctrltitle{GET /competition/\emph{id}/submission}
\ctrlauth{AlgoBOWL Student}
\ctrldesc{Generate the Team Name Change, Input/Output Upload, and Input Download Page}

\ctrltitle{POST /competition/\emph{id}/set\_team\_name}
\ctrlauth{AlgoBOWL Student}
\ctrldesc{Handle form submission to change team name}

Makes sure the team name is valid and less than 20 characters, and sets the
team's name. Redirects back to the inputs page.

\ctrltitle{POST /competition/\emph{id}/upload\_input}
\ctrlauth{AlgoBOWL Student}
\ctrldesc{Handle form submission to upload the group's input}

\begin{enumerate}
    \item Converts input to UNIX file format (if in MS-DOS or Mac OS format)
    \item If they are replacing a previous input:
        \begin{enumerate}
            \item Email all groups who have already submitted outputs and let
                them know they will need to redo their input and upload again.
                The text for this email is shown in Figure
                \ref{emailredo}.\footnote{This is not a catch-all solution for
                preventing outputs that don't correspond to the correct input
                (a student could still download an input, the other group
                changes their input, then the student uploads their output),
                but it's the best we can do without having students put unique
                ID's at the top of their inputs and outputs.}
            \item Delete any outputs which have already been submitted
        \end{enumerate}
    \item Store their input in the database
\end{enumerate}

\begin{figure}[H]
\begin{verbatim}
To: student1@mines.edu, student2@mines.edu, student3@mines.edu
From: algobowl@connect.mines.edu
Reply-To: instructor@mines.edu
Subject: [AlgoBOWL] New Input Uploaded for Group {{Group ID}}

Dear {{Team Name}},

Group {{Group ID}} has just uploaded a new input on AlgoBOWL and your
team had an output submitted already.

Your old output has been deleted. At your convenience, please download
their new input, rerun your program on it, and upload the new output.

Best regards,

The AlgoBOWL System
\end{verbatim}
    \caption{The email to be sent when an output redo is needed}
    \label{emailredo}
\end{figure}

\ctrltitle{GET /competiton/\emph{id}/input/input\_group\_\emph{group\_id}.txt}
\ctrlauth{AlgoBOWL Student or AlgoBOWL Admin}
\ctrldesc{Download input for group \texttt{group\_id}}

If Group \texttt{group\_id} has not uploaded an input yet, this page should
abort with a 404 error.

\ctrltitle{GET /competition/\emph{id}/all\_inputs.zip}
\ctrlauth{AlgoBOWL Student or AlgoBOWL Admin}
\ctrldesc{Generate a ZIP archive of all the inputs}

\ctrltitle{POST /competition/\emph{id}/upload\_output}
\ctrlauth{AlgoBOWL Student}
\ctrldesc{Handle form submission to upload the group's output}

Takes a parameter of the input Group ID the output corresponds to. This should:

\begin{enumerate}
    \item Convert the output to UNIX file format (if in MS-DOS or Mac OS format)
    \item Strip the first line off the output and convert it to a numerical
        type. This is the \texttt{score} attribute. If this process fails, it
        should raise an error to the user.
    \item Store the remainder of the output as the \texttt{data} attribute
    \item Set the \texttt{verification} attribute to \emph{Waiting}
    \item Store their output in the database
\end{enumerate}

\ctrltitle{GET /competition/\emph{id}/verification}
\ctrlauth{AlgoBOWL Student}
\ctrldesc{Generate a page with output download links and a button to accept or
reject each of them}

\ctrltitle{GET /competition/\emph{id}/output/output\_from\_\{from\_id\}\_to\_\{to\_id\}.txt}
\ctrlauth{AlgoBOWL Student with Group ID \texttt{to\_id} or AlgoBOWL Admin or
anyone after open verification starts}
\ctrldesc{Download output from group \texttt{from\_id} to group
\texttt{to\_id}}

\ctrltitle{GET /competition/\emph{id}/all\_outputs\_to\_\{to\_id\}.zip}
\ctrlauth{AlgoBOWL Student with Group ID \texttt{to\_id} or AlgoBOWL Admin or
anyone after open verification starts}
\ctrldesc{Generate a ZIP archive of all the outputs to group \texttt{to\_id}}

\ctrltitle{GET /competition/\emph{id}/evaluations}
\ctrlauth{AlgoBOWL Student}
\ctrldesc{Generate student evaluations page}

\ctrltitle{POST /competition/\emph{id}/evaluations}
\ctrlauth{AlgoBOWL Student}
\ctrldesc{Set AlgoBOWL Evaluations for a certain student}

This controller method should verify the submitted evaluations sum to exactly
100, and raise an error otherwise.

\ctrltitle{GET /setup}
\ctrlauth{AlgoBOWL Admin}
\ctrldesc{Geterate AlgoBOWL Setup Page}

\ctrltitle{POST /setup}
\ctrlauth{AlgoBOWL Admin}
\ctrldesc{Process a new AlgoBOWL setup}

This method should create a new competition with the specified teams and
configuration

\ctrltitle{GET /competition/\emph{id}/verif\_editor}
\ctrlauth{AlgoBOWL Admin}
\ctrldesc{Generate verification editing page for the admin}

\ctrltitle{POST /competition/\emph{id}/verif\_editor}
\ctrlauth{AlgoBOWL Admin}
\ctrldesc{Process a verification edit for the admin}

\ctrltitle{GET /competition/\emph{id}/eval\_admin}
\ctrlauth{AlgoBOWL Admin}
\ctrldesc{View evaluations}

\ctrltitle{GET /competition/\emph{id}/evaluations.csv}
\ctrlauth{AlgoBOWL Admin}
\ctrldesc{Download a CSV file with the evaluations}

\section{Interface \& View Design}

\subsection{Tabs}

For students, the tabs should read:

\texttt{Rankings | Submission | Verification | Evaluations}

For admins, the tabs should read:

\texttt{Rankings | Submission | Verification Editor | View Evaluations | Setup}

For users who are not logged in, they should be able to view the rankings page,
but instead of the tabs there is a message:

\emph{You are not currently enrolled on an AlgoBOWL team. If you think you have
received this message in error, please contact your instructor.}

\subsection{Rankings}

The rankings page should view like FRANKENBOWL: \\
\url{https://algobowl.mines.edu/algobowl/rankings}

Note that the view for this page from the old AlgoBOWL will likely not be
anyhow salvageable, as it is far different from Jack's prototype.

\subsection{Submission}

The submission page should be like the one on the current AlgoBOWL: \\
\url{https://algobowl.mines.edu/algobowl/submission}

It should also include a link to download all inputs as a ZIP file.

For the admin user, do not show the Input Upload or Change Team Name sections.

\subsection{Verification}

The verification screen should look like the one on the current AlgoBOWL: \\
\url{https://algobowl.mines.edu/algobowl/verification}

It should also include a link to download all outputs for that team as a ZIP
file.

\subsection{Evaluation}

The evaluation screen should look like the one on the current AlgoBOWL: \\
\url{https://algobowl.mines.edu/algobowl/evaluations}

\subsection{Verification Editor}

The verification editor screen should present the admin with a matrix of
verifications (text boxes containing A, R, or W) and allow the admin to edit
the verifications and resubmit the page.

\subsection{View Evaluations}

The evaluations page should show a table of all evaluations to the admin as
well as a button to download a CSV file?

\end{document}
